\section{Two Optimizers, One Algorithm}\label{sec:opt}

\S\ref{sec:famous} found structural similarities in the code required to implement
classification and clustering algirithms.
This section does the same, for two optimizers. 

A  (1+1) evolutionary algorithm uses 1 current solution
to generate 1 mutant. If the mutant is ``better'' in some way then it becomes the new current solution used in future mutations.

Examples of (1+1) or
simulated annealing, and local search.
Both have their own mechanisms for escaping from
searching dead-ends. Simulated annealing may  jump to a slightly less optimal solution (especially
early in a run when a ``temperature'' constant is very ``hot''). Local search has a reset-retry
policy where, if no progress is being made, it jumps back to the start and searches in a new
randomly selected method.

Another  difference between them is that simulated annealing always mutates an attribute, selected at random. Local search does the same except, a probability $p$, when it freezes all but one of the attributes, then samples across the range of the remaining
attributes. 

Simulated annealing and local search are often presented as separate algorithms.
But as seen here, both can be coded by calling an underlying \textsf{oneplus1} function,
while specializing a \textsf{mutate} and \textsf{accept} function. 
In the following \textsf{restarts=0} means ``never restart'' and
\textsf{restarts=100} means ``restart if no improvements after 100 mutations''.
\begin{minted}{python}
def sa(d, oracle, restarts=0, m=0.5, budget=1000): 
  n = max(1, int(m * len(d.cols.xs)))
  def accept(e, en, h, b): return en<e or rand()<exp(
                                 (e-en)/(1-h/b+1E-32))
  def mutate(s): yield picks(d, s, n)
  return oneplus1(d, mutate, accept, oracle, budget, restarts)

def ls(d, oracle, restarts=100,p=0.5, tries=20, budget=1000): 
  def accept(e, en, *_): return en < e
  def mutate(s):
    c = choice(d.cols.xs)
    for _ in range(tries if rand()<p else 1):
      s = s[:]
      s[c.at] = pick(c, s[c.at])
      yield s
  return oneplus1(d, mutate, accept,  oracle, budget, restarts)
\end{minted}                  
 XXX how are the accepst and ,itates differenr. 

One housekeeping item. Optimizers need a way to say which row is
better. We use \textsf{disty}, which is the Euclidean distance
from a row to an ideal point computed only on goal columns. From
line~\ref{heaven}, each goal column knows its ideal value: 1 for
``+'' columns, 0 for ``-'' columns. \textsf{disty} just walks the
goal columns:

\begin{minted}{python}
def disty(data, row): 
  return minkowski( abs(norm(col, row[col.at]) - col.heaven)
                    for col in data.cols.ys, p=the.p)
\end{minted}

Lower scores are better. \S\ref{sec:using} expands the role of
\textsf{disty} when active learning needs it. Here a one-line
definition is fine.

\subsection{Mutation via \textsf{pick}; scoring via \textsf{nearest}}\label{spick}

You need three things to optimize: a way to perturb a row, a way
to score the result, and a rule for accepting the change. \IT{}
already has the first two.

\textsf{picks} mutates a row by replacing some random independent
columns with new values drawn from \textsf{pick}, the sampler we
saw in \S\ref{sec:famous} for $k$-means++. For \cls{Num} columns,
\textsf{pick} jiggles the value around a Gaussian centered near
the current value. For \cls{Sym} columns, it samples from the
observed distribution.

\begin{minted}{python}
def picks(data, row, n=1): 
  s = row[:]
  for col in sample(data.cols.xs, min(n, len(data.cols.xs))):
    s[col.at] = pick(col, s[col.at])
  return s
\end{minted}

Scoring is harder. Running the real oracle on each candidate is
the whole reason this paper exists. We substitute a cheap
proxy: find the closest already-labeled row in $x$-space, copy
its $y$ values into the candidate, and call \textsf{disty}.

\begin{minted}{python}
def oracleNearest(data, row): 
  near = nearest(data, row)
  for col in data.cols.ys: row[col.at] = near[col.at]
  return disty(data, row)
\end{minted}

Three lines. This is a 1-NN surrogate. The expensive oracle is
called only when active learning (\S\ref{sec:active}) decides a
genuinely new label is worth paying for.

\subsection{(1+1)-EA, simulated annealing, and local search:
three accept rules}

The optimization loop takes a mutate callback, an accept
callback, and an oracle. It yields each new best candidate as it
finds them. A \textsf{restart} parameter lets the search jump to
a fresh starting point if it stalls.

\begin{minted}{python}
def oneplus1(data, mutate, accept, oracle,
             budget=1000, restart=0):
  """(1+1) search: mutate, score, accept."""
  h, best, best_e = 0, None, 1E32
  s, e, imp = choice(data.rows)[:], 1E32, 0
  while h < budget:
    for sn in mutate(s):
      h += 1
      en = oracle(sn)
      if accept(e, en, h, budget):
        s, e = sn, en
      if en < best_e:
        best, best_e, imp = sn[:], en, h
        yield h, best_e, best
      if restart and h - imp > restart:
        s = choice(data.rows)[:]
        e, imp = 1E32, h
        break
\end{minted}

Notice what is not here. There is no commitment to a search
strategy. There is a loop shape (mutate, score, accept, track
best), and the search strategy is whatever \textsf{accept}
decides.

The (1+1) evolutionary algorithm accepts any candidate no worse
than the current one. Plateaus do not stop it.

\begin{minted}{python}
def acc_ea(e, en, h, budget):
  return en <= e
\end{minted}

Local search accepts only strict improvements. Plateaus stop it.

\begin{minted}{python}
def acc_ls(e, en, h, budget):
  return en < e
\end{minted}

Simulated annealing always accepts improvements, and accepts
worse candidates with probability $\exp(\Delta/T)$, where $T$
cools from 1 to 0 as the budget is spent.

\begin{minted}{python}
def acc_sa(e, en, h, budget):
  if en < e: return True
  T = 1 - h/budget
  return T > 0 and rand() < exp((e-en) / (T+1e-9))
\end{minted}

That is the whole story. Three classical metaheuristics, three
one-line callbacks, surrounded by identical code.

\subsection{Empirical: indistinguishable on MOOT}

If three algorithms differ only by a one-line callback, their
results should be similar. We tested this on MOOT (full corpus
described in \S\ref{sec:using}; for now, treat it as 120-plus
real-world SE optimization tasks scored by \textsf{disty}).

Each task got the same per-trial budget. Each optimizer ran 20
times. We recorded final \textsf{disty}.

For each pair of optimizers on each task, we tested whether
their distributions of final scores were statistically
distinguishable using the Cliff's-delta plus Kolmogorov--Smirnov
test from \S\ref{sec:active}. The test rejects ``different''
when the effect is small and the distributions overlap.

% (REVIEWER NOTE: replace with actual numbers from the
% experimental run. The shape of result we expect, following
% Ganguly~\cite{kashin}, is that EA, SA, and LS are not
% statistically separated on the majority of MOOT tasks.)

On most MOOT tasks, no pair of the three is distinguishable.
Where one wins, the effect is small enough to fall under
Sawilowsky's negligible threshold of $0.35\sigma$
\cite{sawilowsky2009new}.

We also ran a fourth condition: pure random search. We replaced
\textsf{accept} with a function that always returns true and
mutated freely. Random loses. But it loses by a margin smaller
than the metaheuristics literature predicts. On a meaningful
fraction of MOOT tasks, it is statistically tied with the named
methods.

\subsection*{The lesson: random is good enough}

The named-vs-random gap is small. The named-vs-named gaps are
mostly invisible.

A reasonable conclusion is that the loop shape matters and the
accept rule barely does. The metaheuristics literature
distinguishes EA from SA from LS as if they were different
algorithms. They are the same algorithm with a different
one-liner. On real SE data, that one-liner is rarely the thing
that determines whether the search succeeds.

A second conclusion is that the substrate did not need new
abstractions. \textsf{pick} was already there for kpp seeding.
\textsf{nearest} was already there for distance queries.
\textsf{disty} was already there for goal-column scoring. The
new code in this section is \textsf{picks} (six lines),
\textsf{oracleNearest} (three lines), and the search loop
\textsf{oneplus1} (about fifteen lines with restart logic).
Under 25 new lines covered three textbook chapters.

This matches what we saw in \S\ref{sec:famous}. The textbook
treats classification, clustering, and optimization as three
separate fields. The code treats them as three queries against
the same incrementally maintained
\cls{Num}/\cls{Sym}/\cls{Data} distributions.
